\abstract{%
\textbf{Problem.} Cell-free DNA (cfDNA) whole-genome sequencing has become the dominant non-invasive substrate for non-invasive prenatal testing, minimal residual disease monitoring, and early cancer detection, but methylation information remains expensive to access because bisulfite conversion is destructive and biased. \citet{liu2024finaleme} closed part of this gap with FinaleMe, a per-sample two-state hidden Markov model that recovers methylation from fragmentomic features alone. Its binary latent state, point-estimate Baum-Welch fit, per-sample training, susceptibility to prior leakage, lower-information CpG regimes, and binarized downstream tissue deconvolution leave room for a fully probabilistic successor.

\smallskip\textbf{Method.} We present \textbf{HAVI-Methyl} (Hierarchical Amortized Variational Inference for Methylation), a hierarchical Bayesian framework for methylation prediction from cfDNA fragmentomics. The full proposed model uses a population locus prior $\beta^{\mathrm{pop}}_\ell$, a per-sample shift $\delta_s$, and a per-(sample, locus) latent $\beta_{s,\ell}$ on the logit scale, with direct-count, pseudo-count, end-motif, and coverage likelihoods chosen according to the observation regime. Inference is designed as stochastic variational inference with approximate moment-matched natural-gradient updates on the population layer and an amortized normalizing-flow variational posterior $q_\phi(\eta_{s,\ell}\mid F_{s,\ell})$ over per-locus fragment bags, encoded by a permutation-invariant Set Transformer and optional frozen long-range DNA embeddings. Prior-leakage control is handled by a variational information-bottleneck penalty, mQTL anchor losses, and counterfactual augmentation. Calibrated prediction sets come from posterior-predictive checks plus a split-conformal wrapper, and the proposed tissue-of-origin module uses posterior methylation uncertainty rather than binarized point estimates.

\smallskip\textbf{Results.} The completed experiments in this manuscript are synthetic experiments from the released simplified harness, not a completed full real-data benchmark. On 12 simulated samples and 300 loci across coverages $\{0.1\times,1\times,5\times,30\times\}$, the simplified HAVI-Methyl variant achieves Pearson $r$ of $0.333$ / $0.782$ / $0.868$ / $0.964$ versus a FinaleMe-style HMM baseline at $0.162$ / $0.524$ / $0.768$ / $0.961$, with consistent RMSE reductions and stronger recovery on a synthetic low-information locus proxy. In a synthetic prior-leakage stress test, VIB reduces prior attribution from $5.9\%$ to $0.6\%$ of the predicted DMR signal, and adding mQTL anchors reduces it to $0.04\%$. A simplified continuous deconvolution proxy reduces synthetic tissue-fraction RMSE from $0.129$ to $0.007$; full joint Dirichlet-head evaluation remains part of the planned benchmark.

\smallskip\textbf{Contributions.} We provide ELBO, reparameterization-gradient, and approximate natural-gradient derivations; give formal propositions or proof sketches for lower-bound validity, conditional instrumental-anchor identifiability, VIB-bounded prior leakage, conformal marginal coverage, surrogate SVI convergence, hierarchical-pooling variance reduction, and information-theoretic recoverability limits; and introduce a chromatin-aware cfDNA fragmentation-simulator specification. We preserve the planned real-data benchmarking program against the paired WGS/WGBS data of \citet{liu2024finaleme} and public tissue atlases, with ablations, calibration diagnostics, and compute-planning estimates clearly marked as planned work.
}
